% AFT, epub/pagination, 7/10/00
% DXR PRETEX 3/1/00
%%%APT-TXCX(janet)04/26/00
%%%%%%%%%%%%%%%%%%%%%%%%%%%% ADG FINALS 7.5.00 %%%%%%%%%%%
\documentclass[final,oneeqnum]{siamltex}
\usepackage{amsfonts}

%\oddsidemargin=-.35in
%\evensidemargin=-.35in
%\topmargin=.25in
%\textwidth=7.2in
%\textheight=8.3in
%\headsep=.3in
%\advance\textheight by 1.3in
%\advance\topmargin by -.9in
%\def\baselinestretch{1.5}
%\advance\textwidth by 2.2in
%\pagestyle{headings}
%\setcounter{secnumdepth}{-2}

\title{A Monotonicity Property Involving
\protect\boldmath$\,_{3}F_{2}$ and Comparisons of the Classical Approximations
of Elliptical Arc Length\thanks{Received by the editors January 15, 1999;
accepted for publication (in revised form) December 8, 1999; published
electronically July 11, 2000.
\URL sima/32-2/35050.html}}  

\author{Roger W.\ Barnard\thanks{Department of 
Mathematics, Texas Tech University, 
Lubbock, TX  79409 
(barnard@\break math.ttu.edu, pearce@math.ttu.edu)} 
\and Kent Pearce\footnotemark[2]
\and Kendall C.\ Richards\thanks {Department of 
Mathematics, Southwestern University, 
Georgetown, TX 78626 (richards@\break southwestern.edu) }}

\begin{document}

\maketitle
\vspace{-1.35in}
\slugger{sima} {2000}{32}{2}{403--419}
\vspace{1.05in}

\setcounter{page}{403}

\begin{abstract}
Conditions are determined under which
$  \,_{3}F_{2}\left(-n,a,b;a+b+2,\varepsilon -n+1;1\right)$
is a monotone  function of $n$ satisfying
$a b\cdot\,_{3}F_{2}\left(-n,a,b;a+b+2,\varepsilon -n+1;1\right)
 \geq a b\cdot\,_{2}F_{1}\left(a,b;a+b+2;1\right).$ Motivated by
 a conjecture of Vuorinen [{\it Proceedings of Special Functions and
Differential Equations}, K. S. Rao, R. Jagannathan, G. Vanden Berghe, J. Van der Jeugt, eds., 
Allied Publishers, New Delhi, 1998], the corollary that
$\,_{3}F_{2}(-n,-\frac{1}{2},-\frac{1}{2};1,\varepsilon -n+1;1)
\geq \frac{4}{\pi},$  for $1>\epsilon
\geq\frac{1}{4}$ and $n\geq 2,$ is used to determine surprising
hierarchical relationships among the 13 known historical approximations
 of the  arc length of an ellipse.\ This
 complete list of inequalities  compares
  the  Maclaurin series coefficients of  $\,_{2}F_{1}\,$
  with the coefficients of each of the  known approximations,
 for which maximum errors
 can then be established.\ These approximations range over
four centuries from Kepler's in 1609 to Almkvist's in 1985
and include two from Ramanujan.
\end{abstract}

\begin{keywords} 
hypergeometric, approximations, elliptical arc length 
\end{keywords}

\begin{AMS}
33C, 41A
\end{AMS}

\begin{PII}
S003614109935050X
\end{PII}

\pagestyle{myheadings}
\thispagestyle{plain}
\markboth{ROGER W.\ BARNARD, KENT PEARCE, AND KENDALL C.\ RICHARDS}
{A MONOTONICITY PROPERTY INVOLVING $\,_{3}F_{2}$}

\section{Introduction}
\noindent Let $L(x,y)$ be the arc length of an ellipse with
\label{`semiaxes'} semiaxes of  length $x$ and $y$ (with $x \geq  y > 0$)
and let $\displaystyle  \lambda \equiv \frac{x-y}{x+y}.$
In 1742, Maclaurin  \cite{MC} determined that
\begin{eqnarray}  L(x,y) =
\pi
(x+y)\,\cdot\,_{2}F_{1}\left(-\frac{1}{2},-\frac{1}{2};1;\lambda^2\right),
\label{AX}
\end{eqnarray} where  ${\,}_{2}F_{1}$ is the   hypergeometric function
defined by
\[
 {\,}_{2}F_{1}\left(a,b;c;z\right)\equiv 1+\sum_{n=1}^{
\infty }\frac{(a)_{n} (b)_{n}z^{n}}{ (c)_{n}n!}
\]
 with the Appell (or Pochhammer) symbol $(a)_{n}\equiv
a(a+1)\cdots(a+n-1)$ for $n\geq 1$ and $(a)_{0}\equiv 1,$ $a\neq 0$.\  
(For more background  information,
 see \cite{AB}, \cite{R},  \cite{BPR}, and the recent
 survey
article  \cite{BARN} by the first author.)
 In \cite{AB}, Almkvist and Berndt
 compiled and presented  the
 list
 of  the approximations in Table \ref{Approximations} for
 \[
 G(\lambda) \equiv
\,_{2}F_{1}\left(-\frac{1}{2},-\frac{1}{2};1;\lambda^2\right)=\frac{L(x,y)}{
\pi(x+y)}.
\]
%%%%%%%Table 1.1

 \noindent These approximations and their historical and recent connections
to the
approximations of $\pi$ can be found in the Borweins' book \cite{BB}.\ Another
excellent source for historical and current studies of these
topics is the  book \cite{AVV} by Anderson, Vamanamurthy, and Vuorinen.

Recently, several
inequalities between various mean values and the hypergeometric
function were proved in \cite{BB}, \cite{VV}, and the dependence
of the hypergeometric function
$\,_{2}F_{1}(a,b;c;z)$  on its parameters was studied in \cite{ABRVV},
\cite{AQV}.\ These results led to a conjecture of Vuorinen
 (see \cite{MV}) concerning Muir's  approximation $A_{5}$.
 Vuorinen conjectured
(see \cite{MV}) that
\begin{table}
%%%{\scriptsize %%%%%% TOO BIG --ADG
{\tiny
\caption{\rm{Approximations of $G(\lambda)\equiv\,_{2}F_{1}(-\frac{1}{2},-\frac{1}{2};
1;\lambda^2)$ (see [2]).}}
\label{Approximations} 
\begin{tabular}{||l|l|c||} \hline
 \hspace*{.4in  }Discoverer(s)                  &\hspace*{.6in}
Approximation   & {$\delta_{p} =$ first nonzero term} \\
  \hspace*{.1in} and year of discovery&    \hspace*{.9in} $A_{p}(\lambda)$
&  in the Maclaurin series for \\
\hspace*{.0in}    &  & $\Delta_{p}(\lambda) \equiv A_{p}(\lambda)-G(\lambda)$
\\\hline\hline &&\\[-.05in] Kepler, 1609 &
$\displaystyle A_{1}(\lambda) \equiv (1-\lambda^2)^{1/2}$ & $\displaystyle
\delta_{1}=-\frac{3}{4}\lambda^2$\\[-.05in] &&\\\hline
&&\\[-.05in] Euler, 1773 &$\displaystyle A_{2}(\lambda)
\equiv(1+\lambda^2)^{1/2}$& $\displaystyle
\delta_{2}=\frac{1}{4}\lambda^2$\\[-.05in] &&\\\hline &&\\[-.05in] Sipos,
1792&$ \displaystyle A_{3}(\lambda)
\equiv\frac{2}{1+\sqrt{1-\lambda^2}}$ &$\displaystyle
\delta_{3}=\frac{7}{64}\lambda^4$\\[-.1in]
 Ekwall, 1973&&\\\hline &&\\[-.05in] Peano, 1889     & $\displaystyle
A_{4}(\lambda)
\equiv\frac{3}{2}-\frac{1}{2}(1-\lambda^2)^{1/2}$ &$\displaystyle
\delta_{4}=\frac{3}{64}\lambda^4$\\[-.05in] &&\\\hline &&\\[-.05in]
Muir, 1883& $\displaystyle A_{5}(\lambda) \equiv \left(\frac{(1+\lambda)^{3/2}
+(1-\lambda)^{3/2}}{2}\right)^{2/3}
$&$\displaystyle\delta_{5}= -\frac{1}{64}\lambda^{4}$\\[-.05in] &&\\\hline
&&\\[-.05in] Lindner, 1904-1920      &
$\displaystyle A_{6}(\lambda) \equiv \left(1+\frac{\lambda^2}{8}\right)^2
$&$\displaystyle\delta_{6}=
-\frac{1}{2^8}\lambda^{6}$\\[-.05in] Nyvoll, 1978&&\\\hline &&\\[-.05in]
Selmer, 1975     & $\displaystyle A_{7}(\lambda)
\equiv 1+\frac{\lambda^2/4}{1-\lambda^2/16} $&$\displaystyle\delta_{7}=
-\frac{3}{2^{10}}\lambda^{6}$\\[-.05in] &&\\\hline
&&\\[-.05in] Ramanujan, 1914     & $\displaystyle A_{8}(\lambda)
\equiv 3-\sqrt{4-\lambda^2}  $&$\displaystyle
\delta_{8}=-\frac{1}{2^9}\lambda^{6}$\\[-.05in] Fergestad, 1951 &&\\\hline
&&\\[-.05in] Almkvist, 1978 &$\displaystyle
A_{9}(\lambda) \equiv  2\frac{\left(1+\sqrt{1-\lambda^2}\right)^2
+\lambda^2\sqrt{1-\lambda^2}}{\left(1+\sqrt{1-\lambda^2}\right)\left(1+\sqrt[4]{1-\lambda^2}
\right)^2}$&$\displaystyle
\delta_{9}=\frac{15}{2^{14}}\lambda^{8}$\\[-.05in] &&\\\hline &&\\[-.05in]
Bronshtein $\mbox{and}$ Semendyayev, 1964     &
$\displaystyle A_{10}(\lambda) \equiv\frac{64-3\lambda^4}{64-16\lambda^2}
$&$\displaystyle\delta_{10}=
-\frac{9}{2^{14}}\lambda^{8}$\\[-.05in]
 Selmer, 1975&&\\ \hline &&\\[-.05in] Selmer, 1975    & $\displaystyle
A_{11}(\lambda)
\equiv
\frac{3}{2}+\frac{\lambda^2}{8}-\frac{1}{2}\left(1-\frac{\lambda^2}{2}\right)^{1
/2}
$&$\displaystyle\delta_{11}=
-\frac{5}{2^{14}}\lambda^{8}$\\[-.05in] &&\\\hline &&\\[-.05in] Jacobsen
$\mbox{and}$
Waadeland, 1985     & $\displaystyle
A_{12}(\lambda) \equiv \frac{256-48\lambda^2 -21\lambda^4}{256-112\lambda^2
+3\lambda^4}  $&$\displaystyle\delta_{12}=
-\frac{33}{2^{18}}\lambda^{10}$\\[-.05in] &&\\\hline &&\\[-.05in]
 Ramanujan, 1914    & $\displaystyle A_{13}(\lambda) \equiv 1+\frac{3\lambda^2
}{10+\sqrt{4-3\lambda^2}}  $&$\displaystyle
\delta_{13}=-\frac{3}{2^{17}}\lambda^{10}$\\[-.05in] &&\\\hline
\end{tabular}
\\[0.05in]
%\label{Approximations}
%\caption{\rm{(see [2]) Approximations of }
%$G(\lambda)\equiv\,_{2}F_{1}\left(-\frac{1}{2},-\frac{1}{2};1;\lambda^2\right)$.
}
\end{table}
\begin{eqnarray} 
A_{5}(\lambda) \leq G(\lambda)\quad
 \mbox{for all}\,\lambda \,\in \,[0,1] \label{eq1}.
\end{eqnarray} 
That is, Vuorinen conjectured that $A_{5}$
is a {\em lower bound} for $G$.
This conjecture was recently proved by the authors in \cite{BPR}
which has become
 the genesis of
  the present article.\  Moreover, the results here attest to the adage that
a single
conjecture may have many ramifications.
Also, note that $A_5$ is one of the
mean values studied in \cite{VV}.\ More approximations for
hypergeometric functions  in terms of such mean values are actively being
sought.
For example, let $\nu \,\in \,{\bf R}\backslash\{0\}$  and
define 
\[
M_{\nu}(\lambda) \equiv
\left[\frac{(1+\lambda)^\nu +(1-\lambda)^\nu}{2}\right ]^{1/\nu}.
\]
H.\ Alzer \cite{ALZER} originally made  the following conjecture.

{\sc Conjecture}.\ {\em
The inequalities \begin{eqnarray} M_{\alpha}(\lambda) \leq
G(\lambda) \leq
M_{\beta}(\lambda) \ \ \mbox{hold  for all $\lambda \,\in\, (0,1)$}
\label{MALZ}
\end{eqnarray}
if and only if} 
\[
\alpha \leq 3/2\quad \mbox{{\em and}}\quad \beta \geq (\ln 2)/ \left(\ln
\frac{\pi}{2}\right) \approx 1.53.
\]
As noted by Alzer \cite{ALZER}, it follows from our results
(see the set of inequalities in  expression
(\ref{eqzz}))
that  (\ref{MALZ}) holds with
$\alpha = 3/2$ and $\beta =2$.
Moreover,  for a fixed $\lambda,$  $M_{\nu}(\lambda)$ is an increasing function
of $\nu$.\ Thus  it
follows  that (\ref{MALZ}) holds for all $\alpha \leq 3/2$
and $\beta \geq 2$.\ It can be shown that $\alpha = 3/2$ is
sharp.%  while the sharpness of
%$\beta = (\ln 2)/(\ln \frac{\pi}{2})$ was recently verified by Alzer and Qiu \cite{ALZER}.

\section{Main results}

\noindent  In an earlier paper (see \cite{BPR}),  the authors
were able to verify  inequality (\ref{eq1}) by
working with the original version of  Vuorinen's
conjecture   in  terms of the eccentricity
(see (\ref{eqzx}) and (\ref{eqVz})).\ In this direction,
a generating function argument (motivated by \cite{AG}) was used  to
obtain the following general result (which will also  be applied in
this paper to obtain Theorem \ref{thm2}).

\begin{theorem}[see \cite{BPR}]
\label{thm1}%\noindent{\bf Theorem 1.}
Suppose $a,b>0$.\ Then for any
$\varepsilon $ satisfying $1>\varepsilon \geq\frac{ab}{a+b+1},$ it
follows that
\[
\,_{3}F_{2}\left(-n,a,b;a+b+1,\varepsilon -n +1;1\right)\geq 0, 
\] 
for all
integers $n
\geq 1,$ where $\,_{3}F_{2}$ is the generalized hypergeometric function.\end{theorem}

In light of the conjecture in (\ref{eq1}), the following
question naturally arises:
\begin{quote} {\em  Which  of the remaining  approximations given in Table {\rm\ref{Approximations}}
are  upper bounds or  lower bounds  for $G$}?
\end{quote}

 \noindent An attempt to compare an approximation $A_{p}$ with $G$
 motivates  an  analysis of  the term $\delta_{p}$  (the
first nonzero term in the Maclaurin series representation for the error
function
$\Delta_{p}(\lambda)\equiv A_{p}(\lambda)-G(\lambda)$).
 What information does $\delta_{p}$
 provide? Certainly the leading term can be viewed as a measure of
 accuracy of the given approximation, and the error
function $\Delta_{p}(\lambda)$  will have the same sign as
 $\delta_{p}$
 for {\em sufficiently small} $\lambda$.\ For example, $\delta_{1} <0$ and
  it follows directly that $A_{1}$ is   a lower bound for
 $G,$ as  Kepler intended  (see \cite[p. 599]{AB}).\ In this case,
 the sign of
$\delta_{1}$ is indicative of the sign of  $\Delta_{1}(\lambda)$
{\em for all}  $\lambda \, \in \,[0,1]$.\ Almkvist and
Berndt  proved (see \cite[p. 603]{AB}) that Ramanujan's first
estimate $A_{8}$ is a lower bound for $G$ by proving the
significantly
 stronger result that  the nonzero Maclaurin series coefficients of
 $\Delta_{8}$ all have the same (negative) sign.\ A
numerical investigation suggests that a similar  trait might be
shared by other   approximations given in Table \ref{Approximations}.\  In this
article, it will be shown that all of the approximations  given
in Table \ref{Approximations} satisfy the following property:
\begin{quote} {\em The sign of the error function $\Delta_{p}(\lambda)$
coincides with the sign of the leading term
$\delta_{p}$ \underline{for all  $\lambda\, \in \,[0,1]$}.}
\end{quote} Moreover,  for all but two of the approximations, it will
be established
 that the nonzero Maclaurin series coefficients of $\Delta_{p}$ all have the
 same  sign as $\delta_{p}$.\ (Only  Euler's approximation and Muir's
 approximation fail to satisfy this condition.)
 As a consequence of the forthcoming results,
 each  function
$|\Delta_{p}|$ is  a strictly increasing function of
 $\lambda,$ for $p=1,\ldots,13$.\ Therefore,
$0=|\Delta_{p}(0)| < |\Delta_{p}(\lambda)|< |\Delta_{p}(1)|$  for
all $\lambda \, \in \,(0,1)$.
 For example, the maximum error for Ramanujan's second  estimate
 is $|\Delta_{13}(1)| = |\frac{14}{11} - \frac{4}{\pi}|
\approx 0.000512 $ and satisfies
 $|\Delta_{13}(1)|< |\Delta_{p}(1)|$
 for $p=1,\ldots,12$.\  In this direction, we will prove the following
 three propositions.

%\pagebreak
\begin{proposition}
\label{prop1}%\noindent{\bf Proposition 1.}
\ Let $G(\lambda)
\equiv \sum_{n=0}^{\infty}\alpha_{n}\lambda^{2n} $ and
 $A_{p}(\lambda) \equiv 
\sum_{n=0}^{\infty}\beta^{(p)}_{n}\lambda^{2n}$
 where
$\alpha_{n} \equiv   (\frac{(-1/2)_{n}}{n!})^2$ and each
$A_{p}$ is
defined as in Table {\rm\ref{Approximations}}.\ Then
\begin{eqnarray}
 \beta^{(12)}_{n}\leq \alpha_{n} \leq \beta^{(9)}_{n}\quad \mbox{ for all
integers $n
\geq 0$.}  \nonumber
\end{eqnarray} Therefore, the error functions $|\Delta_{9}|$  and
$|\Delta_{12}|$ are
strictly increasing and
\end{proposition}
\[
 A_{12}(\lambda) \leq G(\lambda)\leq A_{9}(\lambda) \quad \mbox{for all
$\lambda$}
 \,\in \,[0,1].
\]

\begin{proposition}
\label{prop2}%\noindent{\bf Proposition 2.}
\ Let $G(\lambda) \equiv
\sum_{n=0}^{\infty}\alpha_{n}\lambda^{2n} $ and
 $A_{p}(\lambda) \equiv
\sum_{n=0}^{\infty}\beta^{(p)}_{n}\lambda^{2n}$
 where
$\alpha_{n} \equiv (\frac{(-1/2)_{n}}{n!})^2$ and each
$A_{p}$ is
defined as in Table {\rm\ref{Approximations}}.\ Then
\begin{eqnarray}
 \beta^{(1)}_{n}\leq \beta^{(6)}_{n}\leq \beta^{(7)}_{n}\leq
\beta^{(8)}_{n}\leq
\beta^{(10)}_{n} \leq \beta^{(11)}_{n}\leq \beta^{(13)}_{n}\leq \alpha_{n}
\leq\beta^{(4)}_{n}\leq\beta^{(3)}_{n}  \nonumber
\end{eqnarray} 
\noindent for all integers $n \geq 0$. 
Therefore, the corresponding error functions $|\Delta_{p}|$ are
strictly increasing and
\[
 A_{1}(\lambda) \leq A_{6}(\lambda) \leq A_{7}(\lambda) \leq A_{8}(\lambda)
 \leq A_{10}(\lambda) \leq A_{11}(\lambda) \leq
A_{13}(\lambda) \leq G(\lambda)\leq  A_{4}(\lambda) \leq  A_{3}(\lambda) 
\]

\noindent {for all $\lambda$}   \, $\in$ \,$[0,1]$.
\end{proposition}

The next proposition addresses the two remaining estimates: Euler's
approximation $A_{2}$ and  Muir's approximation
$A_{5}$.
 The claim will be made that
\begin{eqnarray}
 A_{5}(\lambda) \equiv \left(\frac{(1+\lambda)^{3/2} +
 (1-\lambda)^{3/2}}{2}\right)^{2/3}\leq G(\lambda)\leq
(1+\lambda^2)^{1/2}\equiv A_{2}(\lambda) 
\label{eqzz}
\end{eqnarray} 
\noindent {for all}\ $\lambda$ \, $\in$ \,[0,1].\ As we have noted, the
 nonzero Maclaurin series coefficients of
$\Delta_{2}$ and $\Delta_{5}$ (as functions of $\lambda$) do
not have   constant sign.\   In order to verify the inequalities
in (\ref{eqzz}), we make use of the known fact due to Landen and Ivory
(e.g., see \cite[p. 598]{AB}) that
\begin{eqnarray}
G(\lambda)\equiv\,_{2}F_{1}\left(-\frac{1}{2},-\frac{1}{2};1;\lambda^2\right)
 =\frac{2x}{x+y}
\cdot\,_{2}F_{1}\left(\frac{1}{2},-\frac{1}{2};1;\xi^2\right), \label{eqzx}
\end{eqnarray} 
where  $\lambda \equiv (x-y)/(x+y)$ and
 $\xi\equiv (1/x){\textstyle\sqrt{x^2-y^2}}$
is the \label{`eccentricity'}
eccentricity of the original ellipse (see (\ref{AX})).
 Without loss of generality,  assume that $1=x\geq y\geq 0$.\ A change of
variable
 from $\lambda$ to $\xi$  can be
accomplished in (\ref{eqzz}) by using 
(\ref{eqzx}) and
 the substitutions  $\lambda =(1-y)/(1+y)$ and $y=\textstyle\sqrt{1-\xi^2}$.
 Multiplying through by $(1+y)/2$
 and
 simplifying, we see that the inequalities in (\ref{eqzz}) are equivalent to
\begin{eqnarray}
\left( \frac{1+(1-\xi^2)^{3/4}}{2}\right) ^{2/3} \leq
\,_{2}F_{1}\left(\frac{1}{2},-\frac{1}{2};1;\xi^{2}\right)\leq
 (1-\xi^2/2)^{1/2} 
\label{eqVz}
\end{eqnarray}
\noindent {for all}\ $\xi$ \, $\in$ \,[0,1]. (The first inequality in (\ref{eqVz}) is the original version of Vuorinen's
conjecture \cite{MV}.)

It is interesting to note that one can show that the  functions
in (\ref{eqVz}) can
be shown to satisfy the stated inequalities  by
establishing that the  coefficients of their respective Maclaurin series,
{\em expanded in powers of $\xi$}, satisfy the
corresponding inequality relationships.
\noindent In view of the preceding discussion, we now state the following proposition.
\vfil

\begin{proposition}[see \cite{BPR}]
\label{prop3}%\noindent{\bf Proposition 3.} (see \cite{BPR})
\  Let $G$ and $A_{p}$
be as defined in Table {\rm\ref{Approximations}} and let
\begin{eqnarray}
1+\sum_{n=1}^{\infty}b_{n}\xi^{2n} &\equiv& \left(
\frac{1+(1-\xi^2)^{3/4}}{2}\right) ^{2/3}
\quad \mbox{and} \label{eqRR}\\
 1+\sum_{n=1}^{\infty}c_{n}\xi^{2n} &\equiv& (1-\xi^2/2)^{1/2}.\ \label{eqRR3}
\end{eqnarray} 
It follows that
\begin{eqnarray}
 b_{n}\leq \frac{(1/2)_{n}(-1/2)_{n}}{n!\cdot n!} \leq c_{n} \quad
 \mbox{for all integers $n\geq 1$.}   \nonumber
\end{eqnarray} 
Therefore, $(\ref{eqVz})$ holds and is equivalent to 
$A_{5}(\lambda) \leq G(\lambda)\leq A_{2}(\lambda)\ 
\mbox{for all $\lambda$}   \,\in \,[0,1].\ $\end{proposition}
%\vspace{.1in}

{\it Remark.} If we apply the identity in
(\ref{eqzx}) with $\lambda = (1-\textstyle\sqrt{1-\xi^2})/(1+\textstyle\sqrt{1-\xi^2}),$
 the definition of $A_2,$ and simplify, we obtain
$\Delta_{2}(\lambda) =
2[(1-\xi^2/2)^{1/2}
-\,_{2}F_{1}(\frac{1}{2},-\frac{1}{2};1;\xi^{2})]/(1\break +\textstyle\sqrt{1-\xi^2})$.
Proposition \ref{prop3} implies that $(1-\xi^2/2)^{1/2}
-\,_{2}F_{1}(\frac{1}{2},-\frac{1}{2};1;\xi^{2})$
is a strictly increasing
function of $\displaystyle \xi$.\ Therefore  $\Delta_{2}(\lambda)$ is a
strictly increasing function of $\displaystyle \xi$.
Since $\xi =
\frac{2\sqrt{\lambda}}{1+\lambda}$
is a strictly increasing
function of $\lambda$ on [0,1],  it follows that $|\Delta_{2}|$ is  a strictly
increasing
function of $\lambda$.
A similar argument can be
applied to $|\Delta_{5}|$.

Although some of the inequalities in the above
propositions are straightforward, several proved to be surprisingly
challenging to verify.\  In particular, the effort  involving Almkvist's
approximation
$A_{9}$ precipitated the discovery of  some deeper results involving the
generalized hypergeometric function $\,_{3}F_{2}\,,$
which are also of independent interest.\   In this direction, our main
general results are as follows.
%\vspace{.1in}

\begin{theorem}
\label{thm2}%\noindent{\bf Theorem  2.}
Let  $1>a\geq b>-1$ and
$1>\varepsilon\geq\,\frac{(a+1)(b+2)}{a+b+4}$.\  Then
$\displaystyle T_{n}\equiv\,_{3}F_{2}\left(-n,a,b;a+b+2,\varepsilon
-n+1;1\right)$ satisfies
\end{theorem}
\[ 
a b (T_{n}- T_{n+1})\geq 0  \quad  \mbox{for all integers $n
\geq 2$.} 
\]

\begin{corollary}
\label{cor1}%\noindent{\bf  Corollary   1.}
\ Let  $1>a\geq b>-1$ and
$1>\varepsilon\geq\,\frac{(a+1)(b+2)}{a+b+4}$.\ Then
$\displaystyle T_{n}\equiv \\ \,_{3}F_{2}\left(-n,a,b;a+b+2,\varepsilon
-n+1;1\right)$ satisfies
\end{corollary}
\[
a b T_{n} \geq a b T_{n+1} \geq a b\cdot\,_{2}F_{1}\left(a,b;a+b+2;1\right)
\quad \mbox{for all integers $n
\geq 2$.} 
\]
\begin{corollary}
\label{cor2}%\noindent{\bf  Corollary 2.}
\  Let  $1>\varepsilon\geq\frac{1}{4}$.
Then $
T_{n}\equiv\,_{3}F_{2}\left(-n,-\frac{1}{2},-\frac{1}{2};1,\varepsilon
-n+1;1\right)$ satisfies
\end{corollary}
\[
T_{n} \geq T_{n+1}\geq
\frac{4}{\pi} \quad \mbox{for all integers $n
\geq 2$.} 
\]

\section{Verification of coefficient inequalities} 
\mbox{}

{\it Proof of Proposition {\rm\ref{prop1}}}.\ {\em Part {\rm I:} Almkvist's
Approximation $A_{9}$.} \
Let $s\equiv   (1-\lambda^2)^{1/2}$ and
$\beta_{n}\equiv   \beta^{(9)}_{n}$.\  It follows  that
\begin{eqnarray}
 A_{9}(\lambda) = 2\left[\frac{(1+s) +(1-s)s}{(1+\sqrt{s})^2}\right] =
 \sum_{n=0}^{\infty}\beta_{n}\lambda^{2n},  \nonumber
\end{eqnarray} 
which implies that
\begin{eqnarray}
 2(1+2s-s^2) = (1 + 2\sqrt{s} +
s)\sum_{n=0}^{\infty}\beta_{n}\lambda^{2n}.
\label{eq8}
\end{eqnarray} 
By replacing $s$  by $(1-\lambda^2)^{1/2}$ and applying
$ (1-\lambda^2)^{q} = \sum_{n=0}^{\infty}\frac{(-q)_{n}}{n!}\lambda^{2n}, $
we may change (\ref{eq8}) to the form
\begin{eqnarray} 
2\lambda^2 &+&
4\sum_{n=0}^{\infty}\frac{(-1/2)_{n}}{n!}\lambda^{2n} \nonumber \\
&&=\sum_{n=0}^{\infty}\beta_{n}\lambda^{2n}+
2\sum_{n=0}^{\infty}\sum_{k=0}^{n}\frac{(-1/4)_{n-k}}{(n-k)!}\beta_{k}\lambda^{2
n}+
\sum_{n=0}^{\infty}\sum_{k=0}^{n}\frac{(-1/2)_{n-k}}{(n-k)!}\beta_{k}\lambda
^{2n}.\  \nonumber
\end{eqnarray} 
Equating the coefficients of $\lambda^{2n},$ we obtain
 $\beta_{0}=1,$ $\beta_{1}=1/4,$ and
\begin{eqnarray} 
4\frac{(-1/2)_{n}}{n!}= \beta_{n} +
2\sum_{k=0}^{n}\frac{(-1/4)_{n-k}}{(n-k)!}\beta_{k}
+\sum_{k=0}^{n}\frac{(-1/2)_{n-k}}{(n-k)!}\beta_{k}\quad
\mbox{for $n\geq 2$}.\ \nonumber
\end{eqnarray} 
Solving for $\beta_{n},$ we have the recursive relationship
\begin{eqnarray}
\beta_{n} = \frac{(-1/2)_{n}}{n!} -
\frac{1}{2}\sum_{k=0}^{n-1}\frac{(-1/4)_{n-k}}{(n-k)!}\beta_{k}-
\frac{1}{4}\sum_{k=0}^{n-1}\frac{(-1/2)_{n-k}}{(n-k)!}\beta_{k} \quad
\mbox{for $n\geq 2$}.\ 
\label{eq9}
\end{eqnarray}

\noindent We will use (\ref{eq9}) and induction to
show that 
\begin{eqnarray}
\beta_{n} \geq \alpha_{n} \quad \mbox{for all $n\geq 0$}.\ 
\label{eq10}
\end{eqnarray} 
First note that $\beta_{n} = \alpha_{n}$ for $n=0,1,2$.
Now let $n\geq 2$ and suppose that
$\beta_{k} \geq \alpha_{k}$ for all $k=0,\ldots, n-1$.\ Since the
coefficients of
$\beta_{k}$ in (\ref{eq9}) are all
positive, it follows that
\begin{eqnarray}
\beta_{n} \geq \frac{(-1/2)_{n}}{n!} -
\frac{1}{2}\sum_{k=0}^{n-1}\frac{(-1/4)_{n-k}}{(n-k)!}\alpha_{k}-
\frac{1}{4}\sum_{k=0}^{n-1}\frac{(-1/2)_{n-k}}{(n-k)!}\alpha_{k}.\ \nonumber
\end{eqnarray} 
Thus (\ref{eq10}) will be established if we can verify that
\begin{eqnarray}
\quad \frac{(-1/2)_{n}}{n!} -
\frac{1}{2}\sum_{k=0}^{n-1}\frac{(-1/4)_{n-k}}{(n-k)!}\alpha_{k}-
\frac{1}{4}\sum_{k=0}^{n-1}
\frac{(-1/2)_{n-k}}{(n-k)!}\alpha_{k}\geq \alpha_{n} \quad \mbox{for
$n\geq 2$}.\ 
\label{eq12}
\end{eqnarray}
Next we use the identities
$(c)_{n-k} = \frac{(-1)^{k}(c)_{n}}{(1-c-n)_{k}}$ and
$(1)_{n} = n!$ and add the corresponding $n$th term of each
summation to both sides.\ Then (\ref{eq12}) becomes
\begin{eqnarray}
\quad \frac{(-1/2)_{n}}{n!} -
\frac{(-1/4)_{n}}{2\cdot
n!}\sum_{k=0}^{n}\frac{(-n)_{k}}{(5/4-n)_{k}}\alpha_{k}-
\frac{(-1/2)_{n}}{4\cdot
n!}\sum_{k=0}^{n}\frac{(-n)_{k}}{(3/2-n)_{k}}\alpha_{k}\geq
\frac{\alpha_{n}}{4}.
\label{eq13}
\end{eqnarray} 
Now we apply
$\alpha_{k} \equiv \big(\frac{(-1/2)_{k}}{k!}\big)^2$
 and the definition of  $\,_{3}F_{2},$ then divide both
sides of  (\ref{eq13}) by $\frac{-(-1/2)_{n}}{4\cdot n!},$
 and simplify.\ Then inequality (\ref{eq13})
 becomes
\begin{eqnarray}
 P(n)\cdot\,_{3}&F&_{2}\left(-n,-\frac{1}{2},-\frac{1}{2};1,\frac{5}{4}
-n;1\right) \nonumber \\&+&
\,_{3}F_{2}\left(-n,-\frac{1}{2},-\frac{1}{2};1,\frac{3}{2} -n;1\right)
\geq Q(n), \label{eq16}
\end{eqnarray} 
where $P(n) \equiv 2\frac{(-1/4)_{n}}{(-1/2)_{n}}$
 and $Q(n) \equiv 4-\frac{(-1/2)_{n}}{n!}$.
For $n\geq 2,$ these can be shown to satisfy
\begin{eqnarray} 
P(n) &\leq& P(n+1) \quad \mbox{and} \label{eq16.1}\\ Q(n)
&\geq& Q(n+1).
\label{eq16.2}
\end{eqnarray}  
We first note that inequality (\ref{eq16}) can be
confirmed directly for $n=2,\ldots,6$.\  An application of
Corollary \ref{cor2} (to be proved in the following section), with  the respective
values of
$\varepsilon = 1/4$ and $\varepsilon =
1/2,$
 yields
\begin{eqnarray}
\,_{3}F_{2}\left(-n,-\frac{1}{2},-\frac{1}{2};1,\frac{5}{4} -n;1\right)
&\geq& \frac{4}{\pi}
 \quad \mbox{and} \label{eq16.3}\\
\,_{3}F_{2}\left(-n,-\frac{1}{2},-\frac{1}{2};1,\frac{3}{2} -n;1\right)
&\geq& \frac{4}{\pi}
 \label{eq16.4}
\end{eqnarray}
for all $n\geq 2$.\ From inequalities
(\ref{eq16.1})--(\ref{eq16.4}) with
$n\geq 6,$
it follows  that
\begin{eqnarray}
 P(n)\cdot\,_{3}F_{2}\left(-n,-\frac{1}{2},-\frac{1}{2};1,\frac{5}{4}
-n;1\right) &+&
\,_{3}F_{2}\left(-n,-\frac{1}{2},-\frac{1}{2};1,\frac{3}{2} -n;1\right) \nonumber \\
&\geq&  P(6)\frac{4}{\pi}
 + \frac{4}{\pi} \geq Q(6)
\geq Q(n).\  \nonumber
\end{eqnarray} 
Therefore, inequality (\ref{eq16}) holds for all $n\geq 2$
and hence
$\beta^{(9)}_{n}\equiv \beta_{n}\geq
\alpha_{n}$ for all  $n\geq 0$.\ That is, Almkvist's approximation
satisfies the
property that all of the nonzero Maclaurin
series coefficients of
$\Delta_{9}$ are positive.\ This concludes the proof of Part I of
Proposition \ref{prop1}.

{\it Proof of Proposition} \ref{prop1}.
{\em Part {\rm II:}  Jacobsen $\mbox{and}$ Waadeland's  Approximation
$A_{12}$.} \  Now
we seek to show that the approximation
$A_{12}$ satisfies the property that all of the nonzero Maclaurin series
coefficients of
$\Delta_{12}$ are negative.\  Let $a=3,$ $b=-112,$ $c=256,$ and
$D=\sqrt{b^2-4ac}$.\ It follows that
\[
\frac{1}{au^2 + bu + c}=\frac{2a}{D}\left[\frac{1}{2au + b-D} -\frac{1}{2au
+ b+D}\right] = \sum_{n=0}^{\infty}d_{n}u^n
\quad \mbox{for}  \quad |u| < \left|\frac{D+b}{2a}\right|,
\]  
where
\begin{eqnarray}
d_{n}\equiv
\frac{2a}{D}\left[\frac{(-1)^n(2a)^n}{(b-D)^{n+1}}-\frac{(-1)^n(2a)^n}
{(b+D)^{n+1}}\right]=
 \frac{1}{D}\left(\frac{2a}{D-b}\right)^{n+1}\left[\left(\frac{b-D}{b+D}\right)^
{n+1}-1\right].\ \nonumber
\end{eqnarray} 
It follows that $d_{n} >0$ for all $n\geq 0$ and
\begin{eqnarray} 
A_{12}(\lambda) &\equiv& \frac{256-48\lambda^2
-21\lambda^4}{256-112\lambda^2 +3\lambda^4}\nonumber\\ &=&-7
+\frac{2048 - 832\lambda^2}{256-112\lambda^2 +3\lambda^4}  \nonumber \\ &=&
-7 +(2048 - 832\lambda^2)
\sum_{n=0}^{\infty}d_{n}\lambda^{2n}.\  \nonumber
\end{eqnarray} 
Now let $\beta_{n} \equiv \beta^{(12)}_{n}$.\ Then the nonzero
Maclaurin series
coefficients for $A_{12}$ are given by
$\beta_{0} = 1$ and
\begin{eqnarray}
\beta_{n} = 2048d_{n} - 832d_{n-1} \quad \mbox{for all  $n\geq 1$.}  \nonumber
\end{eqnarray}
Since $(x^{n+1}-1)/(x-1) > x$ for
$x\equiv (b-D)/(b+D)> 1,$ it follows easily that
$(2048d_{n})/(832d_{n-1})>1$
for all $n\geq 1$.\ Thus
\begin{equation}
\beta_{n} >0
\quad \mbox{for all  $n\geq 0$.} \label{eq21}
\end{equation}
Direct calculation reveals that $\beta_{n} = \alpha_{n}$ for $n=0,\ldots,4$.
Also note that
\begin{eqnarray} 
(256 - 112\lambda^2 + 3\lambda^4)
\sum_{n=0}^{\infty}\beta_{n}\lambda^{2n}
= 256 - 48\lambda^2
-21\lambda^4.\  \nonumber
\end{eqnarray} 
Hence
\begin{eqnarray}
\sum_{n=3}^{\infty}(256\beta_{n} - 112\beta_{n-1} +
3\beta_{n-2})\lambda^{2n} = 0.\label{eq22}
\end{eqnarray} 
Thus the coefficients of $\lambda^{2n}$ in (\ref{eq22})
are  zero for all $n\geq 3$.\ Solving for
$\beta_{n}$ and using (\ref{eq21}), we have
\begin{eqnarray}
\beta_{n} = (112\beta_{n-1} - 3\beta_{n-2})/256 < \frac{112}{256}\beta_{n-1}
\quad \mbox{for all  $n\geq 3$.}  \nonumber
\end{eqnarray}
Now suppose that $\beta_{n} \leq \alpha_{n}$ for some integer $n\geq4,$  where
${\textstyle \alpha_{n} \equiv (\frac{(-1/2)_{n}}{n!})^2}$.\ Then
\begin{eqnarray}
\beta_{n+1} < \frac{112}{256}\beta_{n}\leq \frac{112}{256}\alpha_{n}=
\frac{112}{256}\frac{\alpha_{n}}{\alpha_{n+1}}\alpha_{n+1}=
\frac{112}{256}\left(\frac{n+1}{n-\frac{1}{2}}\right)^{2}\alpha_{n+1}
\leq\alpha_{n+1}.\  \nonumber
\end{eqnarray} 
Thus $\beta^{(12)}_{n}\equiv
\beta_{n}\leq \alpha_{n}$ for all integers $n\geq 0$.
This concludes the proof of Part II of Proposition \ref{prop1}.\,
\qquad\endproof%\rule{2.4mm}{2.4mm}

Before proving Proposition \ref{prop2}, we first observe that the
nine approximations involved  have the following respective
Maclaurin series  representations (recursive relationships satisfied
by $\beta^{(13)}_{n}$ and $\beta^{(3)}_{n}$ are
developed in the appendix):
\begin{eqnarray}
  A_{1}(\lambda) &\equiv&(1-\lambda^2)^{1/2}=1 +
\sum_{n=1}^{\infty}\frac{(-1/2)_{n}}{n!}\lambda^{2n}, \label{eqA1}\\
\vspace*{-.2in}\nonumber\\   A_{6}(\lambda)
&\equiv&\left(1+\frac{\lambda^2}{8}\right)^2 = 1 + \frac{\lambda^2}{4}+
\frac{\lambda^4}{64},\label{eqA6}\\
\vspace*{-.2in}\nonumber\\
  A_{7}(\lambda) &\equiv&1+\frac{\lambda^2/4}{1-\lambda^2/16} = 1 +
\frac{\lambda^2}{4} +
\sum_{n=2}^{\infty}\frac{1}{2^{4n-2}}\lambda^{2n},\label{eqA7}\\
\vspace*{-.2in}\nonumber\\
  A_{8}(\lambda) &\equiv&3-\sqrt{4-\lambda^2}=1 + \frac{\lambda^2}{4} -
\sum_{n=2}^{\infty}\frac{(-1/2)_{n}}{n!2^{2n-1}}\lambda^{2n},\label{eqA8}\\
\vspace*{-.2in}\nonumber\\
  A_{10}(\lambda) &\equiv&\frac{64-3\lambda^4}{64-16\lambda^2}=1 +
\frac{\lambda^2}{4} +
\sum_{n=2}^{\infty}\frac{1}{2^{2n+2}}\lambda^{2n},\label{eqA10}\\
\vspace*{-.2in}\nonumber\\
  A_{11}(\lambda)
&\equiv&\frac{3}{2}+\frac{\lambda^2}{8}-\frac{1}{2}\left(1-\frac{\lambda^2}{2}\right)^{1/2}=1 +
\frac{\lambda^2}{4} -
\sum_{n=2}^{\infty}\frac{(-1/2)_{n}}{n!2^{n+1}}\lambda^{2n},\label{eqA11}\\
\vspace*{-.2in}\nonumber\\
  A_{13}(\lambda) &\equiv&1+\frac{3\lambda^2 }{10+\sqrt{4-3\lambda^2}}=1 +
\frac{\lambda^2}{4} +
\sum_{n=2}^{\infty}\beta^{(13)}_{n}\lambda^{2n},\label{eqA13}\\
\vspace*{-.2in}\nonumber\\
  A_{4}(\lambda) &\equiv&\frac{3}{2}-\frac{1}{2}(1-\lambda^2)^{1/2}=1 +
\frac{\lambda^2}{4} -
\frac{1}{2}\sum_{n=2}^{\infty}\frac{(-1/2)_{n}}{n!}\lambda^{2n},\label{eqA4}\\
\vspace*{-.2in}\nonumber\\
  A_{3}(\lambda) &\equiv&\frac{2}{1+\sqrt{1-\lambda^2}}=1 +
\frac{\lambda^2}{4} +
\sum_{n=2}^{\infty}\beta^{(3)}_{n}\lambda^{2n}.\label{eqA3}
\end{eqnarray}

{\it Proof of Proposition {\rm\ref{prop2}}}.\ We seek to establish
the following inequalities  regarding the specified Maclaurin
series coefficients:
\begin{eqnarray}
 \beta^{(1)}_{n}\leq \beta^{(6)}_{n}\leq \beta^{(7)}_{n}\leq
\beta^{(8)}_{n}\leq
\beta^{(10)}_{n} \leq \beta^{(11)}_{n}\leq \beta^{(13)}_{n}\leq
\alpha_{n}\leq\beta^{(4)}_{n}\leq\beta^{(3)}_{n}
 \label{eqBB1}
\end{eqnarray}
for all $n\geq 0$.
Referring to (\ref{eqA1})--(\ref{eqA3}), we note
that the inequalities in (\ref{eqBB1}) are trivial for $n=0$
and $n=1$.\  Thus   we must verify (\ref{eqBB1})
 for all  $n\geq 2$.\ The first two inequalities are immediate while the
 next three inequalities follow directly by induction.\  We
 now proceed to prove the remaining inequalities in
(\ref{eqBB1}).

\noindent$\bullet$ Claim I. {\em
 $\displaystyle \beta^{(11)}_{n} \leq  \beta^{(13)}_{n}\leq \alpha_{n}$
for all $n\geq 2$.}
%\vspace{.15in}

\noindent Let $\beta_{n}\equiv  \beta^{(13)}_{n}$ and
 $\gamma_{n} \equiv  \beta^{(11)}_{n},$ where
 $\beta^{(11)}_{n}  \equiv \frac{-(-1/2)_{n}}{n!2^{n+1}}$  for
$n\geq 2$ (see (\ref{eqA11})) and recall that
$\alpha_{n}\equiv (\frac{(-1/2)_{n}}{n!})^{2}$.
The nonzero Maclaurin series coefficients of Ramanujan's
second estimate $A_{13}$ can be shown to satisfy
 (see the appendix)
 $\beta_{0}=1,$ $\beta_{1}=1/4,$  $\beta_{2}=1/64,$ and
\begin{eqnarray}
\beta_{n}   = \phi_{n-1} - 2^{-5}\beta_{n-1} \quad  \mbox{for all}
\ n\geq 3, \quad \mbox{where} \quad \phi_{n} \equiv -
\frac{(-1/2)_{n}(3/4)^n}{16\cdot n!}.\label{eqR}
\end{eqnarray} 
Applying (\ref{eqR}) twice,  we have
\begin{eqnarray}
\beta_{n} &=&\phi_{n-1} - 2^{-5}\phi_{n-2} + 2^{-10}\beta_{n-2} \quad
\mbox{for all} \ n\geq 4.\ \label{eqR2}
\end{eqnarray} 
Direct calculation reveals that Claim I holds for
$n=2,3, 4$.\ That is,
$\gamma_{n} \leq  \beta_{n} \leq \alpha_{n}$ for $n=2,3, 4$.\ Now let
$n\geq 5$ and suppose that
\begin{eqnarray}
\gamma_{k} \leq  \beta_{k}\leq \alpha_{k} \quad \mbox{for all
$k=2,\ldots,n-1.$} \label{eqih}
\end{eqnarray} 
Then (\ref{eqR2}) and (\ref{eqih}) together imply that
\begin{eqnarray}
\phi_{n-1} &-& 2^{-5}\phi_{n-2} + 2^{-10}\gamma_{n-2} \nonumber \\ &\leq& 
\overbrace{\phi_{n-1} - 2^{-5}\phi_{n-2} + 2^{-10}\beta_{n-2}}^{\beta_{n}} 
\leq \phi_{n-1} - 2^{-5}\phi_{n-2} + 2^{-10}\alpha_{n-2}.\ \label{eqw3}
\end{eqnarray} 
It  can be shown (see the appendix) that
\begin{eqnarray}
\gamma_{n}&\leq& \phi_{n-1} - 2^{-5}\phi_{n-2} + 2^{-10}\gamma_{n-2}
 \quad  \mbox{and} \label{eqw3.1}\\
 \alpha_{n}&\geq& \phi_{n-1} - 2^{-5}\phi_{n-2} + 2^{-10}\alpha_{n-2}
\label{eqw3.2}
\end{eqnarray} 
for all  $n\geq 5$.\ Therefore, using inequalities
 (\ref{eqw3})--(\ref{eqw3.2}) and induction, we have
$\gamma_{n} \leq \beta_{n} \leq \alpha_{n}$ for all $n\geq 2$.
This completes the proof of  Claim I.

\noindent$\bullet$  Claim II. {\em $\displaystyle   \alpha_{n}\leq
\beta^{(4)}_{n} \leq
 \beta^{(3)}_{n}$ for all $n\geq 2$}.
\vspace{.15in}

\noindent If we now apply (\ref{eqA4}),  the first inequality in
Claim II
becomes
\begin{eqnarray}
\alpha_{n}\equiv \left(\frac{(-1/2)_{n}}{n!}\right)^{2}\leq
\frac{-(-1/2)_{n}}{2\cdot n!}\equiv \beta^{(4)}_{n}\quad \mbox{for all $n\geq
2$.} \nonumber
\end{eqnarray} 
This is equivalent to
\begin{eqnarray}
\frac{-2(-1/2)_{n}}{n!}\leq  1\quad \mbox{for all $n\geq 2$}  \nonumber
\end{eqnarray} 
which follows by induction.\ The second inequality in Claim
II involves the
 Maclaurin series coefficients of Sipos and Ekwall's approximation $A_{3}$
 which can be shown to satisfy the following
recursive relationship (see the appendix):
$\beta^{(3)}_{0}=1,$ $\beta^{(3)}_{1}=1/4,$ $\beta^{(3)}_{2}=1/8,$ and
\begin{eqnarray}
\beta^{(3)}_{n} =
-\frac{1}{2}\sum_{k=0}^{n-1}\frac{(-1/2)_{n-k}}{(n-k)!}\beta^{(3)}_{k}
 \quad  \mbox{for all} \ n\geq 2.
\label{eqSE}
\end{eqnarray} 
Note that
\begin{eqnarray}
 -\frac{1}{2}\sum_{k=0}^{n-1}\frac{(-1/2)_{n-k}}{(n-k)!}\beta^{(3)}_{k} =
 \frac{-(-1/2)_{n}}{2\cdot
n!}-\frac{1}{2}\sum_{k=1}^{n-1}\frac{(-1/2)_{n-k}}{(n-k)!}\beta^{(3)}_{k}
  \label{eqDD}
\end{eqnarray} 
\noindent{for all} \ $n \geq$ 2, and
\begin{eqnarray}
\frac{-(-1/2)_{n-k}}{2\cdot(n-k)!}\beta^{(3)}_{k} > 0 \quad
\mbox{for} \ k=1,\ldots,n-1.\ \label{eqD2}
\end{eqnarray} 
Therefore, (\ref{eqSE})--(\ref{eqD2}) together yield
\begin{eqnarray}
\beta^{(4)}_{n}\equiv\frac{-(-1/2)_{n}}{2\cdot n!}\leq  \beta^{(3)}_{n}
\quad \mbox{for all} \ n\geq 2.\ \nonumber
\end{eqnarray}
\noindent This concludes the proof of Claim II and Proposition \ref{prop2}.\,
\qquad\endproof%\rule{2.4mm}{2.4mm}

{\it Remarks on the Proof of Proposition {\rm\ref{prop3}}}.\ From  
(\ref{eqRR3}),
we have that
$\displaystyle c_{n} \equiv \frac{(1/2)^{n}(-1/2)_{n}}{n!}$ for all
$n\geq 1$.
By induction, it can be shown that
\begin{eqnarray}
 \frac{(1/2)_{n}(-1/2)_{n}}{n!\cdot n!} \leq c_{n} \quad \mbox{for all
$n\geq 1$.}  \nonumber
\end{eqnarray} 
In an earlier  paper (see \cite{BPR}), the authors use
 the logarithmic derivative  and  Cauchy products to obtain
 the recursive relationship for $b_{n}$
 (with $b_n$ as defined in (\ref{eqRR})) given by
\begin{equation}
b_{n+1}=\frac{1}{2(n+1)}\left[ \left( \frac{5}{4}n-\frac{1}{2}\right)
b_{n}-\sum\limits_{k=0}^{n-2}(k+1)b_{k+1}\frac{\left( -\frac{1}{4}\right)
_{n-k}}{(n-k)!}\right].\ \label{eq4}
\end{equation}  
Theorem \ref{thm1}, together with (\ref{eq4}), was then used
 (see \cite{BPR}) to
 establish that
\begin{eqnarray}
 b_{n} \leq \frac{(1/2)_{n}(-1/2)_{n}}{n!\cdot n!} \quad
\mbox{for all }  n\geq 1.\  \nonumber \, \qquad\endproof%$\rule{2.4mm}{2.4mm}$}
\end{eqnarray}

\section{Proofs of general results involving \protect\boldmath$\,_{3}F_{2}$}

\noindent We will make use of the following classical
 identities which we include for the reader's convenience
($F\,\equiv{\,}_{3}F_{2}$).

\noindent {\sc Identity 1} \{see \cite[p. 440, eq. (33)]{Pr}\}.
\[ 
\hspace{-15pc} F\left(\rho,a,b;c,\sigma;1\right)-F\left(\rho+1,a,b;c,\sigma+1;1\right) 
\]
\vspace{-2pc}
\[
= \frac{-ab(\sigma-\rho)}{c\sigma(\sigma+1)}
\cdot F\left(\rho+1,a+1,b+1;c+1,\sigma+2;1\right).
\]
\noindent {\sc Identity 2} \{see \cite[p. 59, eq. (3.1.1)]{GR}\}.
\[ 
F(-n,a,b;c,d;1)=\frac{(d-b)_{n}}{(d)_{n}}\cdot F(-n,c-a,b;c,1+b-d-n;1).
\]
\noindent {\sc Identity 3} \{see \cite[p. 440, eq. (26)]{Pr}\}.
\[
\sigma\cdot F\left(\rho,a,b;c,\sigma;1\right) =
\rho\cdot F\left(\rho+1,a,b;c,\sigma+1;1\right) + (\sigma -\rho)\cdot
F\left(\rho,a,b;c,\sigma+1;1\right).
\]
\noindent {\sc Identity 4} \{see \cite[p. 82, eq.\ (14)]{R}\}.
\[ 
(a_{1}-a_{2})\cdot F(a_{1},a_{2},a_{3};b_{1},b_{2};z)=
a_{1}\cdot F(a_{1}+1,a_{2},a_{3};b_{1},b_{2};z)-a_{2}\cdot
F(a_{1},a_{2}+1,a_{3};b_{1},b_{2};z).
\]
\noindent {\sc Identity 5} \{see \cite[p. 440, eq. (30)]{Pr}\}.
\[ 
F\left(\sigma,a,b;c,d;1\right)-F\left(\sigma+1,a,b;c,d;1\right)  =
\frac{-ab}{cd}\cdot
F\left(\sigma+1,a+1,b+1;c+1,d+1;1\right).
\]

{\it Proof of Theorem {\rm\ref{thm2}}}.\ Define
$T_{n}\equiv  F\left(-n,a,b;a+b+2,\varepsilon -n +1;1\right),$ where $F\equiv
\,_{3}F_{2}$.\ Let  $1>a\geq b>-1$ and
$1>\varepsilon\geq\,\frac{(a+1)(b+2)}{a+b+4}$.\ For $n\geq2,$
 it follows that
\begin{eqnarray}  
T_{n+1} - T_{n} &=&
 F\left(-n-1,a,b;a+b+2,\varepsilon
-n;1\right)-F\left(-n,a,b;a+b+2,\varepsilon -n +1;1\right)   \nonumber\\
\ &=& \frac{-a b(\varepsilon+1)}{(\varepsilon-n)
(\varepsilon-n+1)(a+b+2)}F\left(-n,a+1,b+1;a+b+3,\varepsilon -n+2;1\right) \nonumber \\ 
&\ & \quad \{ \mbox{using Identity 1
with $\rho = -n-1,$ $\sigma=\varepsilon -n$}\} \nonumber\\
\ &=& \frac{-a b
(\varepsilon+1)}{(n-\varepsilon)(n-\varepsilon-1)(a+b+2)}
\frac{(\varepsilon-n-b+1)_{n}}{(\varepsilon-n+2)_{n}} \nonumber \\ 
&\ & \quad \times F\left(-n,b+2,b+1;a+b+3,b-\varepsilon ;1\right) \nonumber \\ 
&\ & \quad \{ \mbox{using Identity 2}\}\nonumber\\
\ &=& \frac{-a b
(\varepsilon+1)(b-\varepsilon)_{n}}{(n-\varepsilon)
(n-\varepsilon-1)(a+b+2)(-1-\varepsilon)_{n}(b-\varepsilon)} \nonumber \\ 
&\ & \quad \times [(b+1)F\left(-n,b+2,b+2;a+b+3,b+1-\varepsilon;1\right)\nonumber \\
\  &+&(-\varepsilon
-1)F\left(-n,b+2,b+1;a+b+3,b+1-\varepsilon;1\right)],  \label{eqT4} 
\end{eqnarray} 
where (\ref{eqT4}) follows from Identity 3
 (with $\rho = b+1,$ $\sigma=b-\varepsilon$) and the identity
$(1-\alpha-n)_{n} = (-1)^n(\alpha)_{n}$.

\noindent Identity 4 (with $a_{1} = -n$ and $a_{2} = b+1$)
implies that
\begin{eqnarray}
&F&\left(-n,b+2,b+2;a+b+3,b+1-\varepsilon;1\right) \nonumber \\
&\ & \quad \quad =\frac{1}{b+1}[(n+b+1)F\left(-n,b+2,b+1;a+b+3,b+1-\varepsilon;1 \right) \nonumber \\ 
&\ & \quad \quad + (-n)F\left(-n+1,b+2,b+1;a+b+3,b+1-\varepsilon;1\right)].\  \label{eqT5}
\end{eqnarray} 
Now let $G_{n} = F\left(-n,b+2,b+1;a+b+3,b+1-\varepsilon;1\right)$
 and use (\ref{eqT4}) and (\ref{eqT5}).\ Then we
have that
\begin{eqnarray} 
T_{n+1} - T_{n} &=& \frac{-a b
(\varepsilon+1)(b-\varepsilon)_{n}}{(n-\varepsilon)(n-\varepsilon-
1)(a+b+2)(-1-\varepsilon)_{n}(b-\varepsilon)}\nonumber \\
&\ & \quad \times [(n+b+1)G_{n}
- nG_{n-1} +(-\varepsilon-1)G_{n}]\nonumber \\
\ &=&  \frac{-a b
(\varepsilon+1)(b-\varepsilon)_{n}}{(n-\varepsilon)(n-\varepsilon-1)(a+b+2)(-1
-\varepsilon)_{n}(b-\varepsilon)}\nonumber \\
&\ & \quad \times [n(G_{n}
- G_{n-1})   +(b-\varepsilon)G_{n}].\label{eqT6}
\end{eqnarray}

\noindent Applications of Identity 5 (with  $\sigma=-n$) followed by
Identity 2 yield
\begin{eqnarray} 
G_{n} - G_{n-1} &=& F\left(-n,b+2,b+1;a+b+3,b+1-\varepsilon;1\right) \nonumber \\
&-& F\left(-n+1,b+2,b+1;a+b+3,b+1-\varepsilon;1\right) \nonumber \\
&=& \frac{-(b+2)(b+1)}{(a+b+3)(b+1-\varepsilon)}F\left(-n+1,b+3,b+2;a+b+4,b+2-\varepsilon;1\right)\nonumber \\
&=& \frac{-(b+2)(b+1)}{(a+b+3)(b+1-\varepsilon)}\cdot\frac{(-\varepsilon)_{n-1}}
{(b+2-\varepsilon)_{n-1}}\nonumber \\
&\ & \quad \times F\left(-n+1,a+1,b+2;a+b+4,\varepsilon-n+2;1\right).\  \label{eqT9}
\end{eqnarray}
Identity 2 also implies that
\begin{eqnarray} 
G_{n} = \frac{(-\varepsilon)_{n}}{(b+1-
\varepsilon)_{n}}F\left(-n,a+1,b+1;a+b+3,\varepsilon-n+1;1\right).
\label{eqT10}
\end{eqnarray}

\noindent Combining (\ref{eqT6})--(\ref{eqT10}),  we have
\begin{eqnarray} 
T_{n+1} - T_{n}  = \frac{-a b
(\varepsilon+1)(b-\varepsilon)_{n}}{(n-\varepsilon)(n-\varepsilon-1)(a+b+2)(-1
-\varepsilon)_{n}(b-\varepsilon)}\hspace*{1.5in} 
\nonumber\\ 
\times\left[\frac{-n(b+2)(b+1)(-\varepsilon)_{n-1}}{(a+b+3)(b+1-\varepsilon)(b+2
-\varepsilon)_{n-1}}
F\left(-n+1,a+1,b+2;a+b+4,\varepsilon-n+2 ;1\right) \right.\nonumber \\ \left.
+(b-\varepsilon)\frac{(-\varepsilon)_{n}}{(b+1-\varepsilon)_{n}}F\left(-n,a+1,b+
1;a+b+3,\varepsilon-n+1;1\right)\right].
\label{eqT11}
\end{eqnarray}

\noindent Now make use of
$\frac{(b-\varepsilon)_{n}}{(b+1-\varepsilon)_{n}}=
\frac{(b-\varepsilon)}{(n+ b-\varepsilon)},$ \
$\frac{(-\varepsilon)_{n}}{(-1-\varepsilon)_{n}}=
\frac{(-1-\varepsilon+n)}{(-1-\varepsilon)},$ \
$\frac{(-\varepsilon)_{n-1}}{(-1-\varepsilon)_{n}}=
\frac{1}{(-1-\varepsilon)},$ \ and multiply both sides by
$-ab$.\ Then (\ref{eqT11}) becomes
 \begin{eqnarray} 
\hspace*{-1.1in}a b (T_{n} - T_{n+1})  = \frac{(a
b)^2(\varepsilon+1)}{(n-\varepsilon)(n-\varepsilon-1)(a+b+2)(b-\varepsilon)}
\hspace*{1.5in} \nonumber\\  
\times\left[\frac{-n(b+2)(b+1)(b-\varepsilon)}{(a+b+3)(-1-\varepsilon)(n+
b-\varepsilon)}
F\left(-n+1,a+1,b+2;a+b+4,\varepsilon-n+2 ;1\right) \right.\nonumber\\ 
\left.\ +(b-\varepsilon)\frac{(n-1-\varepsilon)(b-\varepsilon)}{(-1-\varepsilon)(n+b-
\varepsilon)}F\left(-n,a+1,b+1;a+b+3,\varepsilon-n+1;1\right)\right]
\nonumber \\
= \frac{(ab)^2}{(n-\varepsilon)(n-\varepsilon-1)(a+b+2)(n+ b-\varepsilon)}\hspace*{2.0in} 
\nonumber\\  
\times\left[\frac{n(b+2)(b+1)}{(a+b+3)}
F\left(-(n-1),a+1,b+2;a+b+4,\varepsilon-(n-1)+1 ;1\right) 
\right.\nonumber\\ 
\left.\ \frac{}{}+(\varepsilon-b)(n-\varepsilon-1)F\left(-n,a+1,b+1;a+b+3,\varepsilon-n+
1;1\right)\right],
\label{eqT13}
\end{eqnarray} 
where  $n+ b-\varepsilon > n-\varepsilon-1>n-2 \geq 0,$
$n-\varepsilon >0,$ and $\varepsilon-b>
\varepsilon-\frac{(a+1)(b+2)}{a+b+4} \geq 0$.
 Since $1> \varepsilon \geq
\frac{(a+1)(b+2)}{a+b+4}>\frac{(a+1)(b+1)}{a+b+3},$
Theorem \ref{thm1} implies that
\begin{eqnarray}
 F\left(-(n-1),a+1,b+2;a+b+4,\varepsilon-(n-1)+1 ;1\right)
&\geq& 0 \quad \mbox{and}  \nonumber \\
F\left(-n,a+1,b+1;a+b+3,\varepsilon-n+1;1\right) &\geq& 0.\  \nonumber
\end{eqnarray}  
Therefore, (\ref{eqT13}) is the product and sum of
nonnegative quantities and thus
\begin{eqnarray} 
a b (T_{n} - T_{n+1}) \geq 0 \quad \mbox{for all integers }
n\geq 2.\ \nonumber \, \qquad\endproof %\rule{2.4mm}{2.4mm}} \nonumber
\end{eqnarray}

\noindent In order to prove Corollary \ref{cor1}, we will make use of the following
two lemmas.

\begin{lemma}\label{lem1}%\noindent{\bf Lemma 1.}
\ Let $n$ be a positive integer and $0 <
\varepsilon <1$.\ Then
\[ 
\frac{(-n)_{k}}{(\varepsilon-n+1)_{k}} \geq 1 \quad \mbox{for all
$k=0,\dots, n-1$.}
\]
\end{lemma}

{\it Proof of Lemma {\rm\ref{lem1}}}.\ Note that the desired inequality holds at
$k=0$.\ Now
 let $n \geq 2$ and  suppose that
\[ 
\frac{(-n)_{k}}{(\varepsilon-n+1)_{k}} \geq 1 \quad \mbox{for some $k$ with
$0\le k\leq n-2$.}\]
 Then
\[
\frac{(-n)_{k+1}}{(\varepsilon-n+1)_{k+1}}=\frac{(-n)_{k}(-n+k)}{(\varepsilon-
n+1)_{k}(\varepsilon-n+1+k)}\geq
\frac{(-n)_{k}}{(\varepsilon-n+1)_{k}} \geq 1.
\qquad\endproof\]%\mbox{\, \rule{2.4mm}{2.4mm}}\]
%\vspace{.1in}

\begin{lemma}
\label{lem2}%\noindent{\bf Lemma 2.}
\ Define
$\psi_{n}(a,b,c,\varepsilon)\equiv
\frac{(a)_{n}(b)_{n}(-n)_{n}}{n!(c)_{n}(\varepsilon-n+1)_{n}}$.\ Let
$(a,b,c,\varepsilon)$ be in the  domain of $\psi_{n}$ for all $n \geq 2$
with  $\varepsilon <c-a-b$.\ Then
 \[\lim_{n\rightarrow\infty}\psi_{n}(a,b,c,\varepsilon) = 0.\]
\end{lemma}

{\it Proof of Lemma {\rm\ref{lem2}}}.\ Since $(1-c-n)_{n} = (-1)^{n}(c)_{n},$ it
follows that
\[
\psi_{n} = \frac{(a)_{n}(b)_{n}(1)_{n}}{n!(c)_{n}(-\varepsilon)_{n}} =
\frac{\Gamma(a+n)\Gamma(b+n)\Gamma(c)\Gamma(-\varepsilon)}{\Gamma(a)\Gamma(b
)\Gamma(c+n)\Gamma(-\varepsilon+n)}n^{c-a-b-\varepsilon}n^{a+b+\varepsilon-c}.
\]
It is known that  (see \cite[p. 257, eq.\ (6.1.46)]{A})
\[
\lim_{n\rightarrow\infty} \frac{\Gamma(r+n)}{\Gamma(s+n)}n^{s-r}= 1.
\] 
If
$a+b+\varepsilon -c<0,$ then
\[
\lim_{n\rightarrow\infty} \psi_{n} =
\lim_{n\rightarrow\infty}\frac{\Gamma(a+n)\Gamma(b+n)\Gamma(c)\Gamma(-\varepsilon)}
{\Gamma(a)\Gamma(b)\Gamma(c+n)\Gamma(-\varepsilon+n)}n^{c-a-b-\varepsilon}\cdot
\lim_{n\rightarrow\infty}n^{a+b+\varepsilon-c}=
0.
\qquad\endproof
\]
%\mbox{\, \rule{2.4mm}{2.4mm}} \]

{\it Proof of Corollary {\rm\ref{cor1}}}.
Let  $1>a\geq b>-1$ and $1>\varepsilon\geq\,\frac{(a+1)(b+2)}{a+b+4}$ and
define
\[
T_{n}\equiv  \,_{3}F_{2}\left(-n,a,b;a+b+2,\varepsilon -n +1;1\right).
\]
Theorem \ref{thm2} implies that the sequence $\{a b T_{n}\}_{n=2}^{\infty}$
is a monotone (nonincreasing) sequence.\ Now define
\[S_{n} \equiv    1+
\frac{(a)_{n}(b)_{n}(-n)_{n}}{n!(a+b+2)_{n}(\varepsilon -n +1)_{n}}+
\sum_{k=1}^{n-1}\frac{(a)_{k}(b)_{k}}{k!(a+b+2)_{k}}.\]
 Using the definition of $\,_{3}F_{2},$ Lemma \ref{lem1}, and the fact that
$\displaystyle \frac{a b (a)_{k}(b)_{k}}{k!(a+b+2)_{k}} \geq 0$ for
$k=1,\ldots,n-1,$ we obtain
 \begin{eqnarray} 
a b T_{n} =
 a b + \frac{a b (a)_{n}(b)_{n}(-n)_{n}}{n!(a+b+2)_{n}(\varepsilon -n +1)_{n}}+
\sum_{k=1}^{n-1}\frac{a b
(a)_{k}(b)_{k}(-n)_{k}}{k!(a+b+2)_{k}(\varepsilon -n +1)_{k}}
 \geq a b S_{n}  \nonumber
\end{eqnarray}
\noindent {for all $n\geq 2$.} 

\noindent  Applying Lemma \ref{lem2} with $c=a+b+2,$ we have
\[
\lim_{n\rightarrow\infty} S_{n} = \,_{2}F_{1}\left(a,b;a+b+2;1\right).
\]
Since $a b T_{n} \geq a b S_{n}$ for all $n\geq
2,$ it follows that  $\{a b T_{n}\}_{n=2}^{\infty}$ is a bounded monotone
 sequence.\  Thus
\[
a b T_{n}\geq \lim_{n\rightarrow\infty} abT_{n} \geq
\lim_{n\rightarrow\infty} a b S_{n} = a
b\cdot\,_{2}F_{1}\left(a,b;a+b+2;1\right)
\quad \mbox{for all } n\geq 2.\, \qquad\endproof\]%\rule{2.4mm}{2.4mm}}\]

{\it Proof of Corollary {\rm\ref{cor2}}}.\ Choose $a=b=-1/2$ and $1>\varepsilon
\geq 1/4$
 and  define
\[
T_{n}\equiv  \,_{3}F_{2}\left(-n,-\frac{1}{2},-\frac{1}{2};1,\varepsilon
-n +1;1\right).
\]
It is known   that (see
\cite[p. 49]{R})
 \[
\,_{2}F_{1}\left(-\frac{1}{2},-\frac{1}{2};1;1\right) = \frac{4}{\pi}.
\]
Corollary \ref{cor1} implies that
\[
T_{n} \geq  T_{n+1} \geq \frac{4}{\pi} \quad \mbox{for all } n\geq 2.\,\qquad\endproof\]
%\rule{2.4mm}{2.4mm}}\]

\section{Appendix}

\subsection{Recursive relationship for  Maclaurin series
coefficients of Ramanujan's second estimate \protect\boldmath$A_{13}$}
Writing $\beta_{n}\equiv \beta^{(13)}_{n} ,$  we have
\begin{eqnarray} 
3\lambda^2(10-\sqrt{4-3\lambda^2})=(A_{13}(\lambda) -
1)(10^2 -(\sqrt{4-3\lambda^2})^2)=
 (96+3\lambda^2)\sum_{n=1}^{\infty}\beta_{n}\lambda^{2n} \nonumber
\end{eqnarray} 
which implies that
\begin{eqnarray} 
10-2\left(1-\frac{3}{4}\lambda^2\right)^{1/2}= (32
+\lambda^2)\sum_{n=1}^{\infty}\beta_{n}\lambda^{2n-2}.\ \nonumber
\end{eqnarray}
 Applying
$(1-x)^{q} = \sum_{n=0}^{\infty}\frac{(-q)_{n}}{n!}x^{n}$
and simplifying yields
\begin{eqnarray} 8
-2\sum_{n=1}^{\infty}\frac{(-1/2)_{n}(3/4)^n}{n!}\lambda^{2n} = 32\beta_{1}
+
\sum_{n=1}^{\infty}(32\beta_{n+1}+\beta{n})\lambda^{2n}.\  \nonumber
\end{eqnarray} 
Thus  $\beta_{0}=1,$ $\beta_{1}=1/4,$ $\beta_{2}=1/64,$ and
\begin{eqnarray}
\beta_{n+1} =   \frac{-(-1/2)_{n}(3/4)^n}{16\cdot n!}
-\frac{\beta_{n}}{32} \quad \mbox{for all} \ n\geq 1.\ \nonumber
\end{eqnarray} 
Letting $\displaystyle \phi_{n} \equiv   -
\frac{(-1/2)_{n}(3/4)^n}{16\cdot n!},$ we obtain
\begin{eqnarray}
\beta_{n+1} = \phi_{n} - 2^{-5}\beta_{n} \quad \mbox{for all } n\geq 2.\ \nonumber \,\qquad\endproof
%$\rule{2.4mm}{2.4mm}$}   \nonumber
\end{eqnarray}

\subsection{Recursive relationship for  Maclaurin series
coefficients of Sipos $\mbox{and}$ Ekwall's estimate \protect\boldmath$A_{3}$}
Writing $\beta_{n}\equiv \beta^{(3)}_{n} $ and using the Cauchy product, we
have
\begin{eqnarray}
2=A_{3}(\lambda)(1+\sqrt{1-\lambda^2})=\sum_{n=0}^{\infty}\beta_{n}\lambda^{
2n} +
\sum_{n=0}^{\infty}\sum_{k=0}^{n}\frac{(-1/2)_{n-k}}{(n-k)!}\beta_{k}\lambda
^{2n}.\ \nonumber
\end{eqnarray} 
Thus $\beta^{(3)}_{0}=1, \ \beta^{(3)}_{1}=1/4, \
\beta^{(3)}_{2}=1/8, \quad \mbox{and} $
\begin{eqnarray}
\beta^{(3)}_{n} =
\frac{-1}{2}\sum_{k=0}^{n-1}\frac{(-1/2)_{n-k}}{(n-k)!}\beta^{(3)}_{k}
\quad
 \mbox{for all } n\geq 2.\, \nonumber \qquad\endproof%$\rule{2.4mm}{2.4mm}$}  \nonumber
\end{eqnarray}

\subsection{Establishing inequality (\ref{eqw3.1})} Let
$\phi_{n} \equiv  \frac{-(-1/2)_{n}(3/4)^n}{16\cdot n!},$
$\gamma_{n} \equiv  \beta^{(11)}_{n} =
\frac{-(-1/2)_{n}}{n!2^{n+1}},$ and  $n\geq 4$.
Inequality (\ref{eqw3.1}) claims that $\displaystyle
\gamma_{n}\leq \phi_{n-1} - 2^{-5}\phi_{n-2} + 2^{-10}\gamma_{n-2}$.

Direct calculation reveals that the desired inequality holds for
$n=4,\dots,7$.
 Now suppose that $n \geq 7$.\ Since $\gamma_{n-2}>0,$ we have
\begin{eqnarray*} 
\gamma_{n}&-& \phi_{n-1} +
2^{-5}\phi_{n-2} -  2^{-10}\gamma_{n-2} < \gamma_{n}- \phi_{n-1}
+ 2^{-5}\phi_{n-2} \nonumber \\
 &=& \frac{-(-1/2)_{n}}{n!2^{n+1}} +\frac{(-1/2)_{n-1}(3/4)^{n-1}}{16\cdot
(n-1)!}
-2^{-9}\frac{(-1/2)_{n-2}(3/4)^{n-2}}{ (n-2)!} \nonumber \\
 &=& -2^{-9}\frac{(-1/2)_{n-2}(3/4)^{n-2}}{(n-2)!}\cdot\left[
\frac{2^9(4/3)^{n-2}(n-3/2)(n-5/2)}{n(n-1)2^{n+1}} -
\frac{2^5(3/4)(n-5/2)}{(n-1)}+1\right] \nonumber \\
 &=& -2^{-9}\frac{(-1/2)_{n-2}(3/4)^{n-2}}{(n-2)!}\cdot\left[
\frac{2^{n+4}(n-3/2)(n-5/2)}{3^{n-2}n(n-1)} -
\frac{24(n-5/2)}{(n-1)}+1\right].
\end{eqnarray*} 
Since $\frac{(n-5/2)}{(n-1)} \geq \frac{1}{2},$ it follows that
\begin{eqnarray}
\frac{2^{n+4}(n-3/2)(n-5/2)}{3^{n-2}n(n-1)} -
\frac{24(n-5/2)}{(n-1)}+1 \leq
\frac{2^{n+4}(n-3/2)(n-5/2)}{3^{n-2}n(n-1)} - 11.\ \nonumber
\end{eqnarray} 
Thus
\begin{eqnarray} 
\gamma_{n}&-& \phi_{n-1} +
2^{-5}\phi_{n-2} -  2^{-10}\gamma_{n-2}\nonumber \\
 &<& -2^{-5}\cdot9\frac{(-1/2)_{n-2}(3/4)^{n-2}}{(n-2)!}\cdot\left[
\frac{2^{n}(n-3/2)(n-5/2)}{3^{n}n(n-1)} - \frac{11\cdot2^{-4}}{3^2}\right] \nonumber \\
&\leq& \frac{-9(-1/2)_{n-2}(3/4)^{n-2}}{32(n-2)!}\cdot\left[
\left(\frac{2}{3}\right)^{n} - \frac{11}{144}\right]  \nonumber \\
&\leq&
\frac{-9(-1/2)_{n-2}(3/4)^{n-2}}{32(n-2)!}\cdot\left[\left(\frac{2}{3}\right
)^{7} - \frac{11}{144}\right] <0.\ \nonumber
\end{eqnarray} 
Hence the  claim in (\ref{eqw3.1}) is established.\,
\qquad\endproof%$\rule{2.4mm}{2.4mm}$

\subsection{Establishing inequality (\ref{eqw3.2})}  
Let $\textstyle \phi_{n} \equiv - \frac{(-1/2)_{n}(3/4)^n}{16\cdot n!},$
$\textstyle \alpha_{n} \equiv (\frac{(-1/2)_{n}}{n!})^2,$ and
$n\geq 4$.

Inequality (\ref{eqw3.2}) claims that $\displaystyle
\phi_{n} - 2^{-5}\phi_{n-1} + 2^{-10}\alpha_{n-1} \leq \alpha_{n+1}$.\  Note
that
\begin{eqnarray*} 
&\phi&_{n} - 2^{-5}\phi_{n-1} +
2^{-10}\alpha_{n-1} - \alpha_{n+1}\nonumber\\
&=& -\frac{(-1/2)_{n}(3/4)^n}{16\cdot
n!}+2^{-5}\frac{(-1/2)_{n-1}(3/4)^{n-1}}{16\cdot (n-1)!}
 + 2^{-10}\left(\frac{(-1/2)_{n-1}}{(n-1)!}\right)^2 -
\left(\frac{(-1/2)_{n+1}}{(n+1)!}\right)^2  \\
&=& \frac{(-1/2)_{n-1}}{(n-1)!}\left\{
\frac{-(n-3/2)3^n}{n2^{2n+4}} + \frac{3^{n-1}}{2^{2n+7}}
+\frac{(-1/2)_{n-1}}{2^{10}(n-1)!}-
\frac{(n-3/2)(n-1/2)(-1/2)_{n+1}}{n(n+1)\cdot(n+1)!}\right\} \\
&=& \frac{(-1/2)_{n-1}}{(n-1)!}\left\{
\frac{3^{n-1}}{2^{2n+4}}\left[\frac{3(3/2-n)}{n}+ \frac{1}{2^3}\right] +
\frac{(-1/2)_{n-1}}{(n-1)!}\left[
\frac{1}{2^{10}}-\frac{(n-3/2)^2(n-1/2)^2}{n^2(n+1)^2}\right] \right\} \\
&=& \frac{(-1/2)_{n-1}}{(n-1)!}\left\{\frac{U(n)}{V(n)} +1\right\}V(n) ,
\end{eqnarray*} 
where  $$\displaystyle  U(n) \equiv
\frac{3^{n-1}}{2^{2n+4}}\left[\frac{3(3/2-n)}{n}+
\frac{1}{2^3}\right]$$  and 
$$
\displaystyle V(n) \equiv
\frac{(-1/2)_{n-1}}{(n-1)!}\left[\frac{1}{2^{10}}-\frac{(n-3/2)^2(n-1/2)^2}{
n^2(n+1)^2}\right].
$$

\noindent It follows that $V(n) > 0$.\ Now let $W(n) \equiv U(n)/V(n)$.\ Since
$(-1/2)_{n-1} < 0,$ we will be finished if we can
show that $W(n)+1 > 0$ for all $n\geq 4$.\ Direct calculation again yields
$W(4) +1>0$.\ For $n\geq4,$ it is easy to check that
\begin{eqnarray} 
W(n+1) -W(n)  &=&
\frac{
\frac{3^{n}}{2^{2n+6}}\left[\frac{3(1/2-n)}{n+1}+
\frac{1}{2^3}\right]}{
\frac{(-1/2)_{n}}{n!}\left[
\frac{1}{2^{10}}-\frac{(n-1/2)^2(n+1/2)^2}{(n+1)^2(n+2)^2}\right]} \nonumber \\
&-&\frac{
\frac{3^{n-1}}{2^{2n+4}}\left[\frac{3(3/2-n)}{n}+
\frac{1}{2^3}\right]}{
\frac{(-1/2)_{n-1}}{(n-1)!}\left[
\frac{1}{2^{10}}-\frac{(n-3/2)^2(n-1/2)^2}{n^2(n+1)^2}\right]
} \nonumber \\
 &=& \left\{\frac{\frac{3^{n-1}}{2^{2n+4}}}{\frac{(-1/2)_{n-1}}{(n-1)
!}}\right\}
\left\{\frac{3n}{4(n-3/2)}Z(n+1) - Z(n)
\right\}, \label{eq77}
\end{eqnarray} 
where $Z(n) \equiv
[\frac{3(3/2-n)}{n}+ \frac{1}{2^3}]/
        [\frac{1}{2^{10}}-\frac{(n-3/2)^2(n-1/2)^2}{n^2(n+1)^2}].$
 Direct calculation reveals that the expression in
(\ref{eq77}) is nonnegative for $n=4$ and $n=5$.
For $n\geq 6,$ it can be shown by a straightforward calculation that
$0<Z(n+1)\leq Z(n)$.\ Hence
$
\frac{3n}{4(n-3/2)}Z(n+1) - Z(n) \leq Z(n+1) - Z(n) \leq 0$ for all $n\geq 6$.
 Thus $W(n+1) - W(n) \geq 0$  for all $n\geq 4$ since
$(-1/2)_{n-1} < 0$.\ Therefore,
 $W(n) + 1 \geq W(4) + 1 >0$  for all $n\geq 4$.\ This establishes the claim in
(\ref{eqw3.2}).\, \qquad\endproof%$\rule{2.4mm}{2.4mm}$

\subsection*{Acknowledgments}
\noindent
 The authors wish to acknowledge
 the referees for
 their helpful suggestions regarding the revision of
this paper.\ The authors also wish to thank H.\ Alzer for useful correspondence.

\begin{thebibliography}{9}

\bibitem{A}  {\sc M.\ Abramowitz and I.A.\ Stegun}, {\em Handbook of Mathematical
Functions},  Dover, New York, 1972.

\bibitem{AB}  {\sc G.\ Almkvist and B.\ Berndt}, {\em Gauss, Landen, Ramanujan, the
arithmetic-geometric mean, ellipses, $\pi,$ and the
Ladies Diary},  Amer.\ Math.\ Monthly, 95 (1988), pp.\ 585--608.

\bibitem{ALZER} {\sc H.\ Alzer},~% Inequalities for the Complete Elliptic Integrals, preprint, 1998.
{\it private communication}, 1998.

\bibitem{ABRVV} {\sc G.D.\ Anderson, R.W.\ Barnard, K.C.\ Richards, M.K.\ Vamanamurthy,
and M.\ Vuorinen}, {\em Inequalities for zero-balanced hypergeometric functions},
Trans.\ Amer.\ Math.\ Soc., 347 (1995), pp.\ 1713--1723.

\bibitem{AVV} {\sc G.D.\ Anderson, M.K.\ Vamanamurthy, and M.\ Vuorinen}, 
{\em Conformal Invariants, Inequalities, and Quasiconformal Mappings},
 John Wiley \& Sons, New York, 1997.

\bibitem{AQV} {\sc G.D.\ Anderson, S.-L.\ Qiu, and M.\ Vuorinen}, {\em Precise estimates
for differences of the Gaussian hypergeometric function},
J.\ Math.\ Anal.\ Appl., 215 (1997), pp.\ 212--234.

\bibitem{AG}  {\sc R.\ Askey,  G.\ Gasper, and M.\ Ismail}, {\em A positive sum from
summability theory},  J.\  Approx.\ Theory, 13
(1975), pp.\ 413--420.

\bibitem{BARN} {\sc R.W.\ Barnard},
{\em On applications of hypergeometric functions}, J.\ Comput.\ Appl.\ Math.,
105 (1999), pp.\ 1--8.

\bibitem{BPR} {\sc R.W.\ Barnard, K.\ Pearce, and K.C.\ Richards}, {\em An inequality
involving the generalized hypergeometric function and the arc length of an
ellipse},
SIAM J.\ Math.\ Anal., 31(2000), pp. 693--699.

\bibitem{BB} {\sc J.M.\ Borwein and P.B.\ Borwein}, {\em Inequalities
 for compound mean iterations with logarithmic asymptotes}, J.\ Math.\ Anal.
Appl.,
177 (1993), pp.\ 572--582.

\bibitem{GR}  {\sc G.\ Gasper and M.\ Rahman},  {\em Basic Hypergeometric Series},
Cambridge University Press, Cambridge, UK, 1990.

\bibitem{MC}  {\sc C.\ Maclaurin}, {\em A Treatise of Fluxions in Two Books}, Vol.
2, {\rm T.W.\ and T.\ Ruddimans}, Edinburgh, 1742.

\bibitem{Pr}  {\sc A.P.\ Prudnikov, Yu.\ A.\ Brychkov, and O.I.\ Marichev}, {\em
Integrals and Series, Vol. {\rm 3:}  More Special Functions},
Gordon and Breach Science Publishers, New York, 1990.

\bibitem{R}  {\sc E.D.\ Rainville}, {\em Special Functions},  Macmillan, New York, 1960.

\bibitem{VV} {\sc M.K.\ Vamanamurthy and M.\ Vuorinen}, {\em Inequalities for means},
J.\ Math.\ Anal.\ Appl., 183 (1994), pp.\ 155--166.

\bibitem{MV}  {\sc M.\ Vuorinen}, {\em Hypergeometric functions in geometric function
theory}, in Proceedings\ of Special Functions and
Differential Equations, {\rm K.S. Rao, R. Jagannathan, G. Vanden Berghe, and  J. Van der Jeugt}, 
eds., Allied Publishers, New Delhi, 1998.

\end{thebibliography}

%{\small{\sc Department of Mathematics, Texas Tech University, Lubbock,
%Texas, 79409}

%{\em E-mail address}: {\tt barnard@math.ttu.edu }

%\vspace{.1in}

%{\sc Department of Mathematics, Texas Tech University, Lubbock, Texas, 79409}

%{\em E-mail address}: {\tt pearce@math.ttu.edu }

%\vspace{.1in}

%{\sc Department of Mathematics, Southwestern University, Georgetown, Texas,
%78626}

%{\em E-mail address}: {\tt richards@southwestern.edu }
%}

\end{document}
